\paragraph{确定\trGrading{}与符号。}\label{sub:determining-grading-and-signs}
有了前面所做的全部准备，$HF^*(L_0^{\#},L_1^{\#})$ 的定义就相当直截了当了。取定\trRegular{}的\trFloerDatum，并如下定义\trGraded{}的\trFloerCochainGroup：对每个 $k\in\ZZ$，令
\begin{equation}\label{eq:graded-floer-cochains}
CF^k(L_0^{\#},L_1^{\#})=\bigoplus_{i(y)=k}|o(y)|_\KK.
\end{equation}
设 $u\in\scr{M}_Z(y_0,y_1)$ 是\trFloerEquation{}的一个非\trStationary{}解。有一条显然的\trShortExactSequence{}$\RR\to T\scr{M}_Z(y_0,y_1)_u\to T\scr{M}^*_Z(y_0,y_1)_u$，其中 $\RR$ 项对应于\trTranslation{}方向（更确切地说，我们的约定是：正的\trGenerator{}为 $\partial u/\partial s$）。这使得我们能够等同 $T_u\scr{M}_Z(y_0,y_1)_u$ 与 $T_u\scr{M}^*_Z(y_0,y_1)_u$ 的最高次\trExteriorPower，从而按\eqref{eq:disc-dim-orientation} 给出 $\scr{M}^*_Z(y_0,y_1)$ 相对于 $o(y_0)\otimes o(y_1)^{\vee}$ 的一个取定\trOrientation。特别地，若 $u$ 是 $\scr{M}^*_Z(y_0,y_1)$ 的一个\trIsolated{}点，我们就得到一个\trIsomorphism
\begin{equation}\label{eq:relative-orientation-cu}
c_u\colon o(y_1)\longrightarrow o(y_0).
\end{equation}
\trBoundaryOperator{}$\partial$ 的 $(y_0,y_1)$-\trCoefficients{}定义为对所有这些 $c_u$ 取 $\KK$-\trNormalization{}之和：
\begin{equation}\label{eq:boundary-operator-normalization}
\partial^{y_0,y_1}=\sum_u\,|c_u|_\KK\colon\ |o(y_1)|_\KK\longrightarrow|o(y_0)|_\KK.
\end{equation}
这是一种花哨的（免基的）说法，意思是：每个 $u$ 都以 $\pm1$ 作出贡献，正负取决于 $c_u$ 的符号。再看一眼\eqref{eq:disc-dim-orientation} 便知，$\partial$ 的\trDegree{}为 $1$，正如所愿。
要看出为什么\eqref{eq:boundary-operator-normalization} 是一个合理的做法，取一对解 $u_1\in\scr{M}_Z(y_0,\tilde{y})$、$u_2\in\scr{M}_Z(\tilde{y},y_1)$。把 $u_2$ 向右移动足够远之后，便可把它们\trGluing{}起来，得到某个 $u\in\scr{M}_Z(y_0,y_1)$。用\eqref{eq:polygon-compactification} 的\trGromov{}\trCompactification{}的语言说，$u$ 是一维\trModuliSpace{}$\scr{M}^*_Z(y_0,y_1)$ 中的一点，紧靠着\trCompactification{}的边界点 $(u_1,u_2)$。从\trGluingMap{}$g$ 的结构可以看出，它的\trDifferential{}$Dg_{(u_1,u_2)}$ 有如下性质：$(\partial u_1/\partial s,\partial u_2/\partial s)$ 被送到\trTranslation{}切向量 $\partial u/\partial s$，而 $(-\partial u_1/\partial s,\partial u_2/\partial s)$ 则被送到一个指离边界点 $(u_1,u_2)$ 的向量（把 $(u_1,u_2)$ 沿相反方向 $(\partial u_1/\partial s,-\partial u_2/\partial s)$ 移动，相当于把两块拉开，这与让\trGluingLength{}趋于无穷是一回事）。由这一点以及\eqref{eq:orientation-diagram} 可知，经由\trIsomorphism{} \eqref{eq:disc-dim-orientation} 从 $c_{u_1}\circ c_{u_2}\colon o(y_1)\to o(y_0)$ 得到的 $T\scr{M}^*_Z(y_0,y_1)_u$ 的\trOrientation，总是指离\trGromov{}边界的。于是，属于同一个\trConnectedComponent{}的两个边界点以相反的映射 $|c_{u_1}\circ c_{u_2}|_\KK$ 作贡献，这就蕴含 $\partial^2=0$。照常，$HF^*(L_0^{\#},L_1^{\#})$ 定义为 $H(CF^*(L_0^{\#},L_1^{\#}),\partial)$。
现在取一个 $(d{+}1)$ \trPointedDisc{}$S$，带\trBrane\trLabel{}$(L_0^{\#},\ldots,L_d^{\#})$。选好\trRegular{}的\trFloerDatum{}与\trPerturbationDatum。在\trModuliSpace{}$\scr{M}_{\scr{S}}(y_0,\ldots,y_d)$ 的每个\trIsolated{}点 $u$ 处，\trRelative\trOrientation{} \eqref{eq:disc-dim-orientation} 给出一个\trIsomorphism{}$c_u\colon o(y_d)\otimes\cdots\otimes o(y_1)\to o(y_0)$。照\eqref{eq:boundary-operator-normalization} 那样进行，我们用它们来定义一个映射
\begin{equation*}
C\Phi_{\scr{S}}\colon CF^*(L_{d-1}^{\#},L_d^{\#})\otimes\cdots\otimes CF^*(L_0^{\#},L_1^{\#})\to CF^*(L_0^{\#},L_d^{\#}).
\end{equation*}
为证明 $C\Phi_{\scr{S}}$ 是一个\trChainMap，只需考察\trGromov{}\trCompactification{}的边界点的两种可能类型：
\begin{itemize}
\item 第一种情形是 $(u_1,u_2)\in\scr{M}_{\scr{S}}(y_0,\ldots,y_n,\tilde{y},y_{n+2},\ldots,y_d)\times\scr{M}^*_Z(\tilde{y},y_{n+1})$。此时，\trTranslation{}向量 $(0,\partial u_2/\partial s)$ 在 $Dg$ 下的像是 $\scr{M}_{\scr{S}}(y_0,\ldots,y_d)$ 的一个指离边界的切向量。于是，经由\eqref{eq:disc-dim-orientation} 从 $c_{u_1}\circ c_{u_2}$ 产生的该\trModuliSpace{}的局部\trOrientation，与标准的指向外侧的\trOrientation{}相差\eqref{eq:orientation-diagram} 中 \trKoszul{}符号的 $(-1)$ 倍（其中 $\mathrm{index}(D_{\scr{S}_2,u_2})=1$）。
\item 第二种情形是 $(u_1,u_2)\in\scr{M}^*_Z(\tilde{y},y_0)\times\scr{M}_{\scr{S}}(\tilde{y},y_1,\ldots,y_d)$，此时 $(\partial u_1/\partial s,0)$ 的像指向边界，而 \trKoszul{}符号是平凡的。
\end{itemize}
合起来，便得到
\begin{equation}\label{eq:chain-map-boundary-formula}
\sum_n(-1)^*\,C\Phi_{\scr{S}}\bigl(y_d,\ldots,\partial(y_{n+1}),y_n,\ldots,y_1\bigr)
=\partial\bigl(C\Phi_{\scr{S}}(y_d,\ldots,y_1)\bigr),
\end{equation}
其中 $*=i(y_{n+2})+\cdots+i(y_d)$。沿同一思路的一个类似但更简单的论证表明，\trCompositionLaw{} \eqref{eq:composition-law} 成立（不带任何符号）。一如\cref{sub:floer-cohomology}中那样，人们用由 $C\Phi_{\scr{S}}$ 诱导的\trCohomology{}层面的映射 $\Phi_{\scr{S}}$ 来定义 $HF^0(L^{\#},L^{\#})$ 中的\trIdentityMorphism（$d=0$）、\trTriangleProduct（$d=2$）以及\trContinuationMap（$d=1$）。
或许还值得一提的，是基于 $S$ 上\trPerturbationDatum{}的\trFamily{}$(K^s,J^s)$（$s\in[0;1]$）的\trChainHomotopy{}构造。我们考察由此得到的\trParametrization{}的\trModuliSpace，并利用其中的\trIsolated{}点来定义一个映射 $h\colon CF^*(L_{d-1}^{\#},L_d^{\#})\otimes\cdots\otimes CF^*(L_0^{\#},L_1^{\#})\to CF^*(L_0^{\#},L_d^{\#})[-1]$。诸\trModuliSpace{}的局部\trOrientation{}来自\eqref{eq:family-moduli-dimension}，而\trGluing{}方面的考虑则依赖\eqref{eq:family-gluing-diagram}，此时 $\scr{R}=[0;1]$。相对于前面的论证，唯一的差别是：在第二种情形，\trKoszul{}符号约定又引入了一个符号 $(-1)^{\dim(\scr{R})}=-1$。结论便是
\begin{equation}\label{eq:chain-homotopy-formula}
C\Phi^1_{\scr{S}}-C\Phi^0_{\scr{S}}
=\sum_n(-1)^*\,h\bigl(y_d,\ldots,\partial(y_{n+1}),y_n,\ldots,y_1\bigr)
+\partial\bigl(h(y_d,\ldots,y_1)\bigr),
\end{equation}
其中 $C\Phi^k_{\scr{S}}$ 是由 $(K^k,J^k)$ 导出的\trChainMap，而 $*$ 与\eqref{eq:chain-map-boundary-formula} 中的一样。要确立基本的\trWellDefinedness{}性质，需要的尽在于此。
沿着\cref{sub:general-riemann-surfaces} 中的讨论，把该理论推广到其他\trSurface{}$S$ 上并无困难，只要 $\hat{S}$ 的每条边界\trCircle{}都含有奇数个无穷远点即可；而这在适当的意义下导出一个 TQFT。若想容许不具有此性质的\trSurface，就必须要么对\trBrane{}施加合适的\trOrientation{}假设，要么使用\eqref{eq:asymptotic-identification} 中的技巧。作为一个例子，取 $S=\bar{H}$（其端为出向的），带\trBrane\trLabel{}$L^{\#}$。\eqref{eq:disc-even-formula} 说的是：每个\trIsolated{}元素 $u\in\scr{M}_{S}(y)$ 都给出一个\trIsomorphism{}$c_u\colon o(\Lambda^{\#}_{y,0},\Lambda^{\#}_{y,1}\otimes\lambda^{\mathrm{top}}(\Lambda_{y,1}))\to\KK$，而这些可以合并成一个\trDegree{}为 $-n$ 的\trChainMap{}$CF^*(L^{\#},L^{\#}\otimes\lambda^{\mathrm{top}}(TL))\to\KK$。它在\trCohomology{}上诱导的映射恰好就是\eqref{eq:brane-duality-map}。
